Este trabalho teve como foco o desenvolvimento de uma solução tecnológica para triagem inicial de indícios de TDA em crianças, por meio da integração entre rastreamento ocular, métricas do Teste de Desempenho Contínuo, \textit{IoT} e IA. A proposta foi efetivamente desenvolvida em formato de plataforma digital gamificada, com acesso via navegador e voltada ao apoio profissional no pré-diagnóstico de crianças entre 10 e 12 anos.

Os resultados obtidos demonstram que a solução funciona na etapa atual do projeto. O jogo foi implementado, integrado aos módulos de coleta em tempo real e capaz de gerar retorno inicial sobre o desempenho atencional dos participantes a partir do processamento automático dos dados. Também se verificou aderência e alinhamento à ODS 3 da ONU, ao ampliar possibilidades de triagem inicial em contexto supervisionado.

No módulo de rastreamento ocular, observou-se precisão satisfatória para a fase atual, confirmando a viabilidade técnica da abordagem. No componente de \textit{IoT}, o funcionamento foi adequado durante os testes realizados. Entretanto, permanecem pontos de melhoria já identificados: aumento da robustez e da estabilidade da predição ocular em diferentes condições de uso, por meio de estudos de bibliotecas e técnicas de \textit{eye tracking}; aperfeiçoamento ergonômico do \textit{case} físico do dispositivo; e implementação do módulo de IA, cuja consolidação depende da confiabilidade dos dados de rastreamento ocular de entrada.

Como continuidade do estudo, devem ser realizados testes em ambiente controlado com profissionais da saúde e crianças da faixa etária-alvo, além da validação progressiva dos modelos de IA para identificação de padrões e emissão de relatórios detalhados. Conclui-se, portanto, que a proposta desenvolvida apresenta aplicabilidade prática e potencial de contribuição para a identificação precoce de indícios de TDA, ao mesmo tempo em que requer aprimoramentos incrementais para atingir maior maturidade técnica e operacional.